\documentclass[12pt]{article}

\usepackage[margin=1in]{geometry}
\usepackage{amsmath, amssymb}
\usepackage{booktabs}
\usepackage{hyperref}
\usepackage{enumitem}
\usepackage{xcolor}
\usepackage{titlesec}
\usepackage{parskip}

\definecolor{taskblue}{RGB}{31, 97, 141}
\definecolor{notegreen}{RGB}{30, 130, 76}
\definecolor{warnred}{RGB}{176, 58, 46}

\hypersetup{colorlinks=true, linkcolor=taskblue, urlcolor=taskblue}

\title{\textbf{Critical Outbreak Size for Disease Importation into the
United States: A Threshold Analysis Framework}\\[0.5em]
\large Research Pipeline and Student Brief}

\author{Principal Investigator: Jose Herrera-Diestra\\
\small Tarleton State University \quad
\texttt{jherreradiestra@tarleton.edu}}

\date{May 2026}

% ============================================================
\begin{document}
\maketitle

\tableofcontents
\newpage

% ============================================================
\section*{Overview and Motivation}
% ============================================================

\textbf{Context.} The FIFA World Cup 2026 importation risk analysis
(Herrera-Diestra, in preparation) estimates the probability of at least
one imported infectious disease case reaching each US host city, given
observed outbreak sizes in source countries. That model answers the
question: \textit{``Given today's incidence, what is the risk?''}

\textbf{This project asks the inverse question:}

\begin{center}
\fbox{\parbox{0.85\textwidth}{
  \centering\vspace{0.3em}
  \textit{For a given source country and disease, how large does an
  outbreak need to be before at least one imported case is expected
  to arrive in the United States? And, given the current outbreak
  trajectory, \textbf{when} does that threshold get crossed?}
  \vspace{0.3em}
}}
\end{center}

This reframing converts the importation model into an
\textbf{early-warning tool}: instead of retrospectively estimating
risk from observed data, it provides a \emph{critical threshold}
$C_c^*$ that public health authorities can monitor in real time.
When surveillance reports that a country's case count has crossed
$C_c^*$, the probability of at least one US importation has become
non-negligible.

\textbf{Why this matters.} A single imported case of a high-severity
pathogen (e.g., Ebola, Marburg, novel influenza) can trigger enormous
public health responses. Knowing in advance \emph{how large} an
outbreak in country $c$ must be to generate even a 5\% importation
probability allows:

\begin{itemize}
  \item Proportionate pre-positioning of hospital resources at
        high-volume entry airports.
  \item Evidence-based decisions on enhanced traveller screening.
  \item Communication to public health officials about when to escalate
        monitoring vs.\ when risk remains negligible.
\end{itemize}

\textbf{Relationship to the WC paper.} The mathematical framework
here is identical to the WC importation model. All base data
(travel volumes, population, disease incidence) and all code are
inherited from that project. This analysis extends it in a new
direction and is intended as a standalone follow-up paper.

% ============================================================
\section{Mathematical Framework}
% ============================================================

\subsection{The Poisson Importation Model (Inherited)}

The WC analysis models disease importation as a Poisson process.
If $\Lambda_{c,d}$ is the expected number of imported cases of
disease $d$ from country $c$ arriving in the US per month, the
probability of observing at least one importation is:

\begin{equation}
  P(X_{c,d} \geq 1) = 1 - e^{-\Lambda_{c,d}}
  \label{eq:poisson}
\end{equation}

The importation intensity is:

\begin{equation}
  \Lambda_{c,d} = N_c \cdot I_{c,d} \cdot p_d
  = N_c \cdot \frac{\rho_d \cdot C_{c,d}}{P_c} \cdot p_d
  \label{eq:lambda}
\end{equation}

where:
\begin{itemize}
  \item $N_c$ --- monthly travelers from country $c$ to the US
        (COR/I-94 data, June average).
  \item $C_{c,d}$ --- reported cases of disease $d$ in country $c$
        per month.
  \item $P_c$ --- population of country $c$.
  \item $\rho_d$ --- reporting adjustment factor for disease $d$
        (corrects for under-reporting).
  \item $p_d$ --- probability that an infected individual travels
        internationally while infectious.
\end{itemize}

\subsection{Inverting the Model: Critical Outbreak Size}

Setting Equation~\ref{eq:poisson} equal to a threshold
$\theta \in (0,1)$ and solving for $C_{c,d}$:

\begin{equation}
  1 - e^{-\Lambda_{c,d}} = \theta
  \implies
  \Lambda_{c,d} = -\ln(1 - \theta)
\end{equation}

Substituting Equation~\ref{eq:lambda} and solving for the
\textbf{critical case count} $C_{c,d}^*$:

\begin{equation}
  \boxed{
    C_{c,d}^* =
    \frac{-\ln(1 - \theta) \cdot P_c}{N_c \cdot \rho_d \cdot p_d}
  }
  \label{eq:critical}
\end{equation}

\textbf{Interpretation.} $C_{c,d}^*$ is the number of active
monthly cases in country $c$ at which the probability of exporting
at least one case of disease $d$ to the US reaches $\theta$.
It should be computed for multiple threshold levels, e.g.:

\begin{center}
\begin{tabular}{lcc}
\toprule
\textbf{Threshold $\theta$} & \textbf{$-\ln(1-\theta)$} &
\textbf{Interpretation} \\
\midrule
0.05  & 0.051 & ``Emerging concern'' \\
0.10  & 0.105 & ``Low risk'' \\
0.50  & 0.693 & ``Even odds'' \\
0.95  & 2.996 & ``Near-certain importation'' \\
\bottomrule
\end{tabular}
\end{center}

Note that $C_{c,d}^*$ varies dramatically across countries and
diseases. A country with large travel volume to the US (e.g.,
Brazil, Mexico) has a much lower threshold than a country with
minimal travel (e.g., DRC).

\subsection{Temporal Component: Time to Threshold}

If an outbreak is growing, the number of cases at time $t$ (in days)
follows an exponential growth model in the early phase:

\begin{equation}
  C(t) = C_0 \cdot e^{r \cdot t}
  \label{eq:growth}
\end{equation}

where $C_0$ is the observed case count at time $t = 0$ (today) and
$r$ is the epidemic growth rate. The growth rate is related to the
basic reproduction number $R_0$ and the mean serial interval $T_s$:

\begin{equation}
  r \approx \frac{\ln R_0}{T_s}
  \label{eq:growth_rate}
\end{equation}

Setting Equation~\ref{eq:growth} equal to $C_{c,d}^*$ and solving
for $t$:

\begin{equation}
  \boxed{
    t^* = \frac{1}{r} \ln\!\left(\frac{C_{c,d}^*}{C_0}\right)
  }
  \label{eq:time_to_threshold}
\end{equation}

\textbf{Interpretation.} $t^*$ is the number of days from today
until the outbreak is expected to cross the importation risk
threshold, given the current trajectory. If $C_0 > C_{c,d}^*$,
the threshold has already been crossed ($t^* < 0$).

\textbf{Important caveat.} The exponential model is a simplification
valid only for the early epidemic phase (before control measures,
depletion of susceptibles, etc.). For more realism, this can be
replaced by a SEIR model, but the exponential approximation is
sufficient for a first-pass early-warning tool.

\subsection{Aggregate US Importation Risk}

The framework above can be applied to individual US entry airports
by replacing $N_c$ with $N_{c,v}$ (arrivals from country $c$ to
city $v$) using T-100 routing fractions. This allows city-level
threshold curves analogous to the WC analysis.

% ============================================================
\section{Diseases to Include}
% ============================================================

The analysis should cover at minimum the four diseases from the
WC paper, plus Ebola as a high-severity case study that motivates
the project. Parameters are listed below; see
Table~\ref{tab:params} for sources.

\begin{table}[h!]
\centering
\caption{Disease parameters for the threshold analysis.
$\rho_d$ = reporting adjustment; $p_d$ = travel-while-infectious
probability; $R_0$ and $T_s$ are used in the temporal model.
Ebola parameters are for the Bundibugyo strain (May 2026 DRC outbreak).}
\label{tab:params}
\begin{tabular}{lccccl}
\toprule
\textbf{Disease} & $\boldsymbol{\rho_d}$ & $\boldsymbol{p_d}$
  & $\boldsymbol{R_0}$ & $\boldsymbol{T_s}$ (days)
  & \textbf{Incidence source} \\
\midrule
Dengue    & 0.3 & 0.50 & 2.0--6.0 & 20 & WHO monthly \\
Malaria   & 0.2 & 0.30 & 1.1--3.0 & 30 & WHO annual \\
Measles   & 0.6 & 0.05 & 12--18   & 12 & WHO annual \\
Pertussis & 0.3 & 0.70 & 5.5--17  & 14 & WHO annual \\
Ebola     & 0.8 & 0.01 & 1.5--2.5 & 15 & WHO situation reports \\
\midrule
\multicolumn{6}{l}{\textit{Optional extensions (Phase 2):}} \\
Mpox      & 0.5 & 0.30 & 1.5--3.0 & 9  & WHO/CDC \\
Cholera   & 0.3 & 0.20 & 1.5--4.0 & 5  & WHO \\
\bottomrule
\end{tabular}
\end{table}

\textbf{Note on $R_0$ uncertainty.} The $R_0$ estimates above are
literature ranges. The temporal analysis should include a sensitivity
analysis varying $R_0$ within its plausible range.

% ============================================================
\section{Data Requirements}
% ============================================================

Most data are already available from the WC project. New data
needed are marked with $\star$.

\subsection{Already Available (WC Project)}

\begin{itemize}
  \item \textbf{Country populations:} \texttt{Data/population2020.csv}
        (2020 baseline; acceptable for threshold calculations).
  \item \textbf{Monthly COR arrivals by country:}
        \texttt{Data/Monthly\_Arrivals\_Country\_of\_Residence\_COR\_1.csv}
        — use mean June arrivals as $N_c$.
  \item \textbf{T-100 routing fractions:}
        \texttt{Data/t100\_routing\_fractions.csv}
        — for city-level extension.
  \item \textbf{Dengue incidence (monthly, country-level):}
        \texttt{Data/dengue-global-data-2025-12-10.xlsx}
  \item \textbf{Malaria incidence (annual, per 1,000):}
        \texttt{Data/Malaria\_National\_Unit\_data.csv}
  \item \textbf{Measles incidence (annual, per 1,000,000):}
        \texttt{Data/Measles reported cases and incidence *.xlsx}
  \item \textbf{Pertussis incidence (annual, per 1,000,000):}
        \texttt{Data/Pertussis reported cases and incidence *.xlsx}
  \item \textbf{Country--region mapping:}
        derived from COR data in \texttt{importationRisk\_main.R}.
\end{itemize}

\subsection{New Data to Obtain ($\star$)}

\begin{enumerate}
  \item[$\star$] \textbf{Ebola incidence data.}
        Download WHO Disease Outbreak News (DON) and situation
        reports for all Ebola outbreaks since 2000.\\
        URL: \url{https://www.who.int/emergencies/disease-outbreak-news}\\
        Format needed: country, date, confirmed cases, deaths.
        For ongoing outbreaks, use the CDC situation summary:\\
        \url{https://www.cdc.gov/ebola/situation-summary/index.html}

  \item[$\star$] \textbf{$R_0$ and serial interval estimates (literature).}
        Compile a table of $R_0$ and $T_s$ estimates from the
        published literature for each disease. Suggested sources:
        \begin{itemize}
          \item Biggerstaff et al.\ (2014) for influenza.
          \item Althaus (2014) for Ebola ($R_0$ West Africa).
          \item WHO Technical Reports for other diseases.
        \end{itemize}

  \item[$\star$] \textbf{Air passenger volumes by country (most recent COR).}
        The existing COR file covers through mid-2025. If a more
        recent extract is available from CBP
        (\url{https://www.cbp.gov/newsroom/stats/cbp-travel-stats}),
        download and add to \texttt{Data/}.

  \item[$\star$] \textbf{Optional: historical importation events.}
        Compile a validation dataset of documented disease
        importations into the US from 2010--2025 (source country,
        disease, date). Use ProMED
        (\url{https://promedmail.org}) and CDC importation
        reports. This will be used in Section~\ref{sec:validation}.
\end{enumerate}

% ============================================================
\section{Task List}
% ============================================================

Tasks are numbered and grouped by phase. Each task has a suggested
time estimate for a graduate student working part-time.

\subsection{Phase 1 --- Analytical Framework (Weeks 1--2)}

\begin{enumerate}[label=\textbf{T\arabic*.}]

  \item \textbf{Set up the project structure.}
    Fork or clone the WC repository
    (\url{https://github.com/jdiestra1977/FIFA_worldCup_2026_risk}).
    Create a new folder \texttt{OutbreakThreshold/} with a
    dedicated R script \texttt{threshold\_analysis.R}.
    Load all shared data files from \texttt{Data/} --- do not
    duplicate them.
    \textit{[1--2 hours]}

  \item \textbf{Implement the critical threshold formula.}
    Code Equation~\ref{eq:critical} as a function in R:
    \begin{verbatim}
    compute_threshold <- function(N_c, P_c, rho_d, p_d,
                                  theta = c(0.05, 0.10, 0.50, 0.95)) {
      -log(1 - theta) * P_c / (N_c * rho_d * p_d)
    }
    \end{verbatim}
    Apply this to all combinations of countries (from COR data)
    and diseases (from Table~\ref{tab:params}).
    Output: a tidy data frame with columns
    \texttt{Country, disease, theta, C\_star}.
    \textit{[3--4 hours]}

  \item \textbf{Compute thresholds for all country $\times$ disease pairs.}
    Run \textbf{T2} for all countries in the COR dataset and all
    five diseases. Use the mean annual COR arrivals (not June-only)
    for $N_c$ here, since this analysis is not WC-specific.
    Identify the top 20 countries per disease with the \emph{lowest}
    $C^*$ (i.e., most at-risk source countries).
    \textit{[2--3 hours]}

  \item \textbf{Sanity check and validation against known events.}
    \label{sec:validation}
    For three well-documented US importation events (e.g.,
    2014 Ebola in Dallas from Liberia; 2019 measles from Israel;
    2022 Ebola from DRC/Uganda), check whether the observed outbreak
    size at time of travel exceeded $C^*$ for $\theta = 0.05$.
    If the model is well-calibrated, the answer should be ``yes''
    for most cases.
    \textit{[4--6 hours; requires historical data from T$\star$4]}

\end{enumerate}

\subsection{Phase 2 --- Temporal Model (Weeks 3--4)}

\begin{enumerate}[label=\textbf{T\arabic*.}, resume]

  \item \textbf{Implement the epidemic growth model.}
    Code Equations~\ref{eq:growth}--\ref{eq:time_to_threshold}.
    For each disease, use the central $R_0$ estimate from
    Table~\ref{tab:params} to compute $r$. Implement a sensitivity
    wrapper that reruns the calculation for the low and high $R_0$
    bounds.
    \textit{[3--4 hours]}

  \item \textbf{Apply to the current DRC/Ebola outbreak (case study).}
    Use the May 2026 CDC situation summary as $C_0$:
    DRC has 347 total suspected/confirmed cases at the time of
    writing. Compute $t^*$ for $\theta \in \{0.05, 0.50\}$.
    Compare to dengue in Brazil as a contrast.
    \textit{[2 hours]}

  \item \textbf{Generate threshold curves over time.}
    For the top 10 highest-risk country $\times$ disease pairs
    (identified in \textbf{T3}), plot $P(\geq 1)$ as a function of
    outbreak size $C$ (x-axis) for a range of 0 to $2 \times C^*$.
    This produces ``risk ramp'' curves that visualise how quickly
    importation probability grows with outbreak size.
    \textit{[3 hours]}

\end{enumerate}

\subsection{Phase 3 --- Figures and Outputs (Weeks 5--6)}

\begin{enumerate}[label=\textbf{T\arabic*.}, resume]

  \item \textbf{Figure 1: Global threshold map.}
    Choropleth world map where the fill colour for each country is
    $\log_{10}(C^*/\text{current incidence})$ for dengue at
    $\theta = 0.50$.
    Ratio $< 1$ (country already above threshold) shown in red;
    ratio $1$--$10$ in orange; $> 10$ in green.
    Use \texttt{ggplot2} + \texttt{sf} + country polygons from
    \texttt{rnaturalearth}.
    \textit{[4--6 hours]}

  \item \textbf{Figure 2: Threshold ranking table (top 20 per disease).}
    Faceted bar chart: x-axis = country, y-axis = $C^*$ (log scale),
    faceted by disease. Annotate current reported incidence as a
    point on each bar to show how close each country is to the
    threshold.
    \textit{[3--4 hours]}

  \item \textbf{Figure 3: Time-to-threshold curves.}
    For the DRC/Ebola and Brazil/Dengue case studies, plot
    outbreak size $C(t)$ vs.\ time (x-axis), with a horizontal
    dashed line at $C^*$ for each $\theta$ level. Shade the
    region where $C(t) > C^*$ in red.
    \textit{[2--3 hours]}

  \item \textbf{Figure 4: Sensitivity analysis.}
    For a representative country $\times$ disease pair, show how
    $C^*$ changes when $\rho_d$, $p_d$, and $N_c$ are each varied
    $\pm 50\%$ independently (tornado plot or similar).
    \textit{[2--3 hours]}

  \item \textbf{Figure 5: City-level extension (optional).}
    Using T-100 routing fractions, compute city-specific thresholds
    $C_{c,v}^*$ for the top 5 US entry cities and the top 5 source
    countries per disease. Display as a heatmap analogous to the
    WC paper.
    \textit{[4--5 hours]}

\end{enumerate}

\subsection{Phase 4 --- Write-up (Weeks 7--8)}

\begin{enumerate}[label=\textbf{T\arabic*.}, resume]

  \item \textbf{Draft manuscript.}
    Target journal: \textit{PLOS Neglected Tropical Diseases} or
    \textit{Emerging Infectious Diseases}. Suggested structure:
    Introduction, Methods (model + data), Results (Figures 1--4),
    Discussion, Supplementary tables (full $C^*$ table for all
    countries $\times$ diseases). Target length: 3,500--4,500 words.
    \textit{[15--20 hours]}

  \item \textbf{Produce supplementary $C^*$ table.}
    A CSV/Excel file with columns:
    \texttt{Country, disease, theta\_05, theta\_10, theta\_50, theta\_95,
    current\_incidence, ratio\_to\_threshold}.
    This is a practical deliverable that can be shared with
    public health departments independently of the paper.
    \textit{[2 hours]}

\end{enumerate}

% ============================================================
\section{Suggested R Script Structure}
% ============================================================

The student should create a single self-contained R script
\texttt{OutbreakThreshold/threshold\_analysis.R} with the following
sections. This mirrors the structure of \texttt{importationRisk\_main.R}
in the parent project.

\begin{verbatim}
# 0. Packages
# 1. Working directory and shared data loading
#    (population, COR arrivals, T-100 routing, disease incidence)
# 2. Disease parameters table
#    (rho, p, R0, T_s for each disease)
# 3. compute_threshold() function  [Eq. 3]
# 4. compute_time_to_threshold() function  [Eq. 5]
# 5. Critical thresholds: all country x disease combinations
# 6. Temporal model: DRC/Ebola and Brazil/Dengue case studies
# 7. Figure 1: Global threshold map
# 8. Figure 2: Threshold ranking (top 20 per disease)
# 9. Figure 3: Time-to-threshold curves
# 10. Figure 4: Sensitivity analysis
# 11. Figure 5: City-level extension (optional)
# 12. Export supplementary table
\end{verbatim}

% ============================================================
\section{Connections to the WC Paper}
% ============================================================

The student should be aware of and refer to the parent analysis
throughout. Key shared components:

\begin{itemize}
  \item The Poisson framework (Eq.~\ref{eq:poisson}--\ref{eq:lambda})
        is identical.
  \item The parameter values in Table~\ref{tab:params} must match
        Table~1 of the WC paper to ensure consistency.
  \item The T-100 routing fractions and COR data are used as-is.
  \item The country name harmonisation table in
        \texttt{t100\_routing\_prep.R} should be reused to avoid
        join failures.
  \item The WC paper should be cited as the methodological origin
        of the importation model.
\end{itemize}

The threshold analysis is \textbf{not} limited to WC-qualified
countries or to June travel. It uses full annual travel volumes and
covers all 180+ countries in the COR dataset, making it a more
general tool.

% ============================================================
\section{Expected Timeline Summary}
% ============================================================

\begin{center}
\begin{tabular}{clc}
\toprule
\textbf{Week(s)} & \textbf{Milestone} & \textbf{Deliverable} \\
\midrule
1--2 & Analytical framework coded and validated & R script §0--4 \\
3--4 & Temporal model + case studies            & R script §5--6 \\
5--6 & All figures generated                    & Figures 1--5 \\
7    & Supplementary table complete             & \texttt{.csv} file \\
8    & First draft manuscript                   & \texttt{.tex} draft \\
\bottomrule
\end{tabular}
\end{center}

\textbf{Total estimated effort:} 60--80 hours for a graduate student
with intermediate R skills and basic epidemiology background.

% ============================================================
\section{Questions to Discuss with the PI}
% ============================================================

Before starting, the student should clarify the following with
Prof.\ Herrera-Diestra:

\begin{enumerate}
  \item Which diseases to prioritise in Phase~1 (all five, or start
        with dengue and Ebola as contrasting examples)?
  \item Whether to use mean annual COR arrivals or June-specific
        arrivals for $N_c$ (the former is more general; the latter
        is more conservative for WC-period applications).
  \item Whether the optional city-level extension (Figure~5,
        Task~T12) is in scope for the first submission.
  \item Target journal and co-authorship expectations.
  \item Access to the parent GitHub repository.
\end{enumerate}

% ============================================================
\end{document}
